\section{Controlled Evidence Across Writer and Provenance Channels}
\label{sec:intervention}
The controlled evidence separates two channels that are often conflated. Historical L1 experiments change an outcome-conditioned memory writer and therefore identify a writer-mode bundle. Our main L2 experiments instead hold retrieval membership, order, and actionable memory bytes fixed and change only whether an explicit truthful source-outcome indicator is masked or revealed to the executor. The masked arm can still contain implicit outcome clues in the memory text; the treatment therefore measures incremental field exposure beyond content, not provenance-versus-no-provenance.

\subsection{Historical L1 bridge}
Two preregistered WebArena bridges motivate the cleaner intervention but do not identify provenance alone. A terminal-success experiment admits four information-equivalent task pairs and completes 48 requests; the success-writer minus failure-writer effect is $+0.1667$ with $p=0.0785$. With only four independent tasks, the smallest exact one-sided sign-flip value is $1/16=0.0625$, so this endpoint is structurally non-decisive under the frozen 0.05 rule. A disjoint early-action experiment admits 23 task pairs and completes 276 requests; its reference-action-match effect is $-0.0942$ with counterevidence-side $p=0.0792$. The two estimands have opposite descriptive signs, and deterministic audits show systematic writer-style differences. We therefore make no endpoint-robust directional L1 claim and never upgrade L1 to provenance-only L2. Full admission rules, pair construction, runtime repair boundaries, tables, and randomization contracts are retained in Appendix~A.2--A.9.

\subsection{Fresh L2 provenance-exposure experiment}
\label{sec:l2fresh}
The main experiment is a new prospective treatment-outcome object and does not reuse historical R4/R6 outcomes or the stopped R5/R19 prefixes. Its source-bank capacity is nevertheless pilot-informed: an earlier 128-source support stop and zero-outcome diagnostic exposed coverage sensitivity, after which we chose the entire frozen 350-task training universe rather than tuning a smaller cutoff. The R53 lineage then rebuilds all 350 memories from zero under Qwen2.5-7B-Instruct and MemRL, one durable memory at a time. Fresh target selection excludes every historical validation signature and occurs before any new validation outcome. The source run completes 350/350 units with 176 success-derived and 174 failure-derived memories. This balance is descriptive; source outcomes are never used to choose target units.

Before any validation environment is reset, we apply native MemRL retrieval ($k=10$, cosine threshold $0.5$) to fresh exact skill-signature clusters from the frozen validation split. All 40 historical primary/utilization signatures are excluded. Of 108 remaining clusters, 106 have at least one eligible retrieved memory with nonempty content, a source task in the full source universe, and a boolean source outcome. Frozen hash ranking selects the first 32 eligible clusters as the primary sample and the next eight as a disjoint utilization sample. No target outcome is observed during this support step.

\paragraph{Utilization first stage.}
The eight utilization clusters each receive five frozen surfaces: no memory (U0), true retrieved memory (U1), null content (U2), reversed content (U3), and shuffled other-task memory (U4). Promotion uses only the first executable action; terminal success is diagnostic. U1 produces a first-action change relative to both U0 and U2 in 5/8 units, while U2 differs from U0 in 3/8. This satisfies the preregistered rule $5\geq3$ and $5\geq3+1$, rejecting the trivial explanation that the entire memory surface is inert before the primary contrast is opened. It is an anti-inertness guard, not proof that the executor semantically interprets the later source-outcome field.

\paragraph{Primary A/B contrast.}
For each of the 32 fresh primary clusters, the frozen retrieval is rendered into two executor views. In A, the explicit source-outcome field is masked while the retrieved content remains visible. In B, the same rows, order, and actionable content additionally reveal the truthful boolean \texttt{source\_outcome\_success} for each row. Similarity, Q estimates, scores, roles, memory IDs, and source task IDs are hidden from both arms; retrieval is not rerun between arms. Arm order is frozen per cluster before treatment. Each arm constructs a fresh task/session, calls \texttt{task.reset(session)} before inference, and OSInteraction's reset constructs a new Docker container with \texttt{remove=True}; \texttt{task.release()} runs in a \texttt{finally} path and terminates any remaining container. Thus task-visible filesystem/user/process state is arm-local. We hold the reset initial environment and feedback mechanism fixed, not realized downstream observations after actions diverge. All 64 Qwen arm runs complete without execution failure.

\begin{table}[t]
\centering
\caption{Fresh L2 terminal-success result. The inference unit is the paired dependency cluster.}
\label{tab:l2fresh}
\small
\begin{tabular}{lrr}
\toprule
Quantity & Outcome field masked (A) & Truthful outcome field revealed (B) \\
\midrule
Terminal successes & 15/32 & 16/32 \\
Success rate & 0.46875 & 0.50000 \\
\midrule
Paired $B-A$ effect & \multicolumn{2}{c}{$+0.03125$} \\
B-only / A-only successes & \multicolumn{2}{c}{1 / 0} \\
Exact two-sided sign test & \multicolumn{2}{c}{$p=1.0$} \\
Paired bootstrap 95\% interval & \multicolumn{2}{c}{$[0,0.09375]$} \\
Preregistered relevance floor & \multicolumn{2}{c}{$|\Delta|\geq0.15$; not met} \\
\bottomrule
\end{tabular}
\end{table}

Thirty-one of 32 pairs have identical terminal outcomes. The only discordant pair favors B, but one discordant unit supplies essentially no directional resolution under the exact two-sided paired test; $p=1.0$ is not affirmative evidence for zero effect. The preregistered bootstrap interval is reported as specified; because the empirical paired-effect distribution contains 31 zeros and one $+1$, we do not interpret its upper endpoint as an equivalence bound. A post-confirmatory conservative paired-risk-difference audit separately gives $[-0.128,0.183]$, so a $+0.15$ population effect is not excluded. The supported statement is therefore descriptive and threshold-based: the observed $+3.125$-point effect did not reach the preregistered 15-point threshold on the frozen 32-pair sample.

\paragraph{Local policy sensitivity and a real Task~252 example.}
A complete-run diagnostic shows that A and B differ in first executable action for 9/32 clusters and in total step count for 7/32, while terminal success differs for only 1/32. This establishes local sensitivity to the field-exposure treatment, but it does not distinguish semantic use of source outcome from generic structured-prompt sensitivity. The binary terminal endpoint is also many-to-one: behaviorally different trajectories can reconverge after later OS feedback and receive the same success label.

The sole Qwen terminal-discordant cluster is OSInteraction Task~252: \emph{create user \texttt{appuser}, create \texttt{/home/appuser/config.cfg} with \texttt{DEBUG\_MODE=true}, and set mode 600}. Its frozen retrieval contains five memories in the same order in both arms: source tasks 74, 403, 208, and 210 are success-derived configuration/user/permission examples; source task 450 is a failure-derived reflection whose visible text already recommends checking environment capabilities, errors, permissions, and ownership. Thus A is not provenance-information-free. Qwen A takes 16 steps and fails; Qwen B, with only the explicit truthful outcome fields added, takes 8 steps and succeeds after adapting to missing \texttt{sudo} and interactive-\texttt{vi} problems. Under Llama, the same Task~252 A/B surfaces yield different first actions but both succeed in 5 steps. We use this outcome-selected case only to illustrate one possible closed-loop pathway; it is not evidence that semantic provenance reasoning explains the aggregate treatment effect.

\subsection{Executor-backbone replication}
\label{sec:l2replication}
To test whether the low-terminal-value pattern is specific to the Qwen executor, we freeze a second prospective experiment before any Llama validation outcome. It reuses the same 350-memory source bank, R54 frozen retrieval membership/content/order, eight utilization clusters, 32 primary clusters, exact-information renderer, terminal evaluator, randomization rule, and complete-only analysis. Importantly, the source memories themselves were written once by the Qwen/MemRL source process and are reused unchanged; this is therefore an executor-backbone replication, not an independent writer- or retriever-backbone replication. The only scientific factor changed at evaluation is the executor, from Qwen2.5-7B-Instruct to Meta-Llama-3.1-8B-Instruct. A source-side native parser qualification passes 3/3 before validation exposure.

The Llama utilization first stage also passes rather than borrowing the Qwen qualification: true memory produces the required memory-specific first-action divergence in 8/8 clusters, while null-memory versus no-memory differs in 5/8, satisfying $8\geq3$ and $8\geq5+1$. Thus a low terminal A/B contrast cannot be dismissed as complete memory-channel inertness, although the placebo divergence also cautions that first-action sensitivity alone does not prove semantic provenance use. All 64 Llama primary arms complete exactly once.

\begin{table}[t]
\centering
\caption{Exact-information L2 results by executor backbone. Models are adjudicated separately; no cross-model pooling is performed.}
\label{tab:l2backbones}
\scriptsize
\begin{tabular}{lrrp{0.34\linewidth}}
\toprule
Executor & $B-A$ & Discordant & 95\% intervals: preregistered bootstrap; conservative paired RD \\
\midrule
Qwen2.5-7B-Instruct & $+0.03125$ & 1/32 & $[0,0.09375]$; $[-0.128,0.183]$ \\
Meta-Llama-3.1-8B-Instruct & $0.00000$ & 4/32 & $[-0.125,0.125]$; $[-0.225,0.225]$ \\
\bottomrule
\end{tabular}
\end{table}

For Llama, two clusters are B-only successes and two are A-only successes, so the observed paired terminal effect is exactly zero and the discordances are directionally balanced. The exact sign test remains $p=1.0$ for both models (omitted from the table for space), but sparse discordance makes that value low-resolution rather than affirmative null evidence. Table~\ref{tab:l2backbones} therefore reports both the frozen bootstrap and the post-confirmatory conservative paired-risk-difference audit. Neither observed point estimate reaches the preregistered $|\Delta|\geq0.15$ threshold, yet neither conservative interval establishes $\pm0.15$ equivalence or excludes a $+0.15$ population effect. The model-specific results are not combined. The replication is correspondingly narrow: local action sensitivity with sparse terminal discordance recurs under a second executor on the same source-bank/retrieval/task/renderer substrate.

\paragraph{L3 transport.} A source-faithful test still requires the motivating financial writer/checkpoint, original memory interface, fixed acquisition content, and paired source runtime; WebArena or MemRL is not a substitute. Appendix~\ref{app:l3} gives the transport contract. L3 therefore remains replication debt.

The replicated L2 result sharpens a limited governance implication. An explicit source-outcome field can alter local action selection, but the observed terminal estimates on these frozen samples do not reach the prespecified 15-point threshold and do not establish an automatic positive trust value. Section~\ref{sec:psmg} therefore treats provenance as auditable control-plane information whose decision weight must be independently justified, rather than as a monotone SUCCESS/FAILURE trust rule. Neither R56 nor the Llama replication tests PSMG efficacy.
